% ================================================================
%  conteudo/introducao.tex
% ================================================================

\section{Introdução}

A quantificação de biomassa microbiana é uma etapa fundamental no monitoramento e controle de bioprocessos industriais, sendo essencial para a compreensão da cinética de crescimento e da produtividade dos cultivos. Entre os métodos mais empregados, a determinação da massa seca e a turbidimetria por densidade óptica (DO) se destacam pela praticidade e confiabilidade. A correlação entre essas duas abordagens, expressa por meio de uma curva padrão, permite a estimativa rápida e não destrutiva da concentração celular durante o processo fermentativo (SCHMIDELL et al., 2001).

\textit{Saccharomyces cerevisiae} é uma das leveduras mais amplamente estudadas e utilizadas na indústria biotecnológica, sendo empregada na produção de etanol, pão, cervejas, vinhos e compostos de interesse farmacêutico. Seu metabolismo versátil, capaz de operar tanto em condições aeróbias quanto anaeróbias, torna essa espécie um modelo clássico para estudos de fisiologia microbiana e desenvolvimento de bioprocessos (TOPALOGLU et al., 2023; WALKER; WALKER, 2018).

O crescimento microbiano pode ser monitorado por diferentes metodologias. O método gravimétrico da massa seca é considerado padrão-ouro por medir diretamente a quantidade de material celular presente na amostra. No entanto, trata-se de um procedimento demorado e que consome volume de amostra. A turbidimetria, por sua vez, baseia-se na relação entre a concentração celular e a dispersão de luz causada pelas células em suspensão, medida como absorbância a 600~nm, sendo um método rápido, simples e amplamente validado para cultivos de leveduras (MIRA et al., 2022).

A elaboração de uma curva padrão que correlacione os valores de absorbância com as respectivas concentrações de biomassa seca permite a conversão direta das leituras espectrofotométricas em dados quantitativos, eliminando a necessidade de determinações gravimétricas a cada ponto amostrado. Essa ferramenta é especialmente útil em acompanhamentos cinéticos, onde a frequência de amostragem é elevada (SCHMIDELL et al., 2001).
